\subsection{Pro-$p$ extensions of the full tame quotient}

The tame finite-generation statement below follows by the standard
curve and gluing methods; the uniform extension consequence F9b
uses the coefficient estimate proved in F8. Both arguments and the
SGA~1 source passed an independent root audit. No historical
priority assertion is made.

\begin{lemma}[F9: a uniform generator bound for the tame quotient; Known method, proved here]
Keep the notation of F8, and let $\pi_1^t(U;X,D)$ denote the quotient
classifying finite covers tamely ramified along the fixed SNC divisor
$D$ on $X$. Let $\Delta$ be the connected dual graph of $D$. Then
\[
 d\bigl(\pi_1^t(U;X,D)\bigr)
 \leq d_X:=\sum_i(2g_i+\nu_i+1)+b_1(\Delta).
 \tag{F9.1}
\]
In particular this is a bound for the full tame group, not only its
maximal prime-to-$p$ quotient.
\end{lemma}
\begin{proof}
Let $N$ be prime to $p$, and use the root stack $\mathfrak X_N$
and reduced exceptional stack $\mathfrak D_N$ of F8. We prove
that every finite quotient of $\pi_1(\mathfrak X_N)$ satisfies
(F9.1), independently of $N$.

\emph{The closed fibre surjects on fundamental groups.}
Let $\mathfrak Y\to\mathfrak X_N$ be a connected finite \'etale
cover. Its coarse space $Y$ is a finite scheme over $X$.
Indeed, locally $\mathfrak X_N=[V/H]$ with $V$ finite over an
affine open $\Spec R\subset X$ and $H=\mu_N$ or $\mu_N^2$.
The pullback of the cover is $W\to V$ finite \'etale, with
$H$-equivariant descent, and its coarse space is
$\Spec\Gamma(W,\mathcal O_W)^H$. This invariant ring is finite
over $R$: it is an $R$-submodule of a finite $R$-module and $R$
is noetherian. These finite coarse schemes glue. Thus $Y$ is
proper over $A$, and it is connected because a coarse-space map
identifies the underlying topological spaces.

For a proper scheme over a henselian local ring, connectedness
is equivalent to connectedness of the closed fibre. For example,
apply the full faithfulness in
\cite[Lemma 58.9.1, tag 0A48]{StacksProperEtaleSpecialization}
to the scheme and its trivial two-sheeted cover: sections of that
cover correspond exactly to open-and-closed decompositions.
Consequently $Y_k$ is connected. Locally in the above quotient
charts, taking invariants commutes with the residue-field base
change by averaging, since $|H|$ is invertible. Therefore $Y_k$
is the coarse space of $\mathfrak Y_k$, which is connected.
Replacing the closed fibre by its reduced support does not change
its underlying topological space. It follows that the restriction
of every connected finite \'etale cover to $\mathfrak D_N$ stays
connected. This is precisely the surjectivity of
\[
 \pi_1(\mathfrak D_N)\longrightarrow\pi_1(\mathfrak X_N).
 \tag{F9.2}
\]
This argument requires only this surjectivity; it does not assume
a full stack version of the proper specialization equivalence.

\emph{A bound for each normalization component.}
The component $\mathfrak C_i$ is the banded $\mu_N$-gerbe over
the root curve $\mathcal C_i$ obtained from $C_i$ by taking roots
at its $\nu_i$ crossing points. In every finite quotient $G$ of
$\pi_1(\mathfrak C_i)$, the image $I$ of the generic relative
$\mu_N$ is cyclic and central. To see centrality, the band supplies
natural automorphisms of every fibre functor; naturality makes
them commute with its automorphisms. Quotienting the corresponding
torsor by $I$ kills the relative gerbe inertia, so the torsor
descends to a finite \'etale $G/I$-torsor on $\mathcal C_i$.
This descent can be checked on a neutralizing \'etale cover of the
gerbe, where it is exactly descent of an equivariant finite \'etale
algebra with trivial $\mu_N$-action.

Every finite \'etale cover of $\mathcal C_i$ restricts to a cover
of $C_i$ minus its crossing points, tamely ramified at those
points. The local chart is $[\Spec k[[z]]/\mu_N]$ with coarse
parameter $z^N$, so the inertia order divides $N$. A connected
cover remains connected on this dense open subset because the
root curve is normal. Hence
$\pi_1^t(C_i\setminus\{\text{crossings}\})$ surjects onto
$\pi_1(\mathcal C_i)$.
SGA~1, Expos\'e~XIII, Corollaire~2.12 gives $2g_i+\nu_i$
generators for this full tame curve group
\cite[original p.~392; retyped pp.~290--292]{SGA1TameCurves}.
It follows that every such $G$ needs at most
\[
                              2g_i+\nu_i+1               \tag{F9.3}
\]
generators: lift generators of $G/I$ and adjoin one generator of
$I$. The compactness criterion for generators gives the same bound
for $\pi_1(\mathfrak C_i)$ itself.

\emph{Gluing the components.}
Finite \'etale covers of $\mathfrak D_N$ are obtained by gluing
finite \'etale covers of its normalization components, with
identifications over the two preimages of each node gerbe.
The local assertion follows on the node chart
\[
 [\Spec k[[u,v]]/(uv)\,/\,\mu_N^2].
\]
The strictly henselian nodal local ring and its two strictly
henselian branches have only trivial finite \'etale covers. Thus
on the quotient chart all three cover categories are finite
$\mu_N^2$-sets, and the gluing statement is the identification
of these sets at the two branches. The assertion descends from
these charts and the smooth locus. Equivalently, it is the usual
normalization gluing of finite \'etale algebras, with equivariant
descent along a finite tame atlas.

For completeness, the resulting generator bound does not require
injectivity of any edge-group map. Choose a spanning tree of
$\Delta$, base points on the components, and paths to the incident
node points. The gluing identifications along the tree allow one
to use a single fibre. The monodromy is then generated by the
images of the component fundamental groups and one gluing
element for every edge outside the tree. If the subgroup they
generated in a finite connected torsor were proper, its orbits
would give a nontrivial decomposition of all component covers
compatible with every gluing identification, hence a decomposition
of that torsor. This contradicts connectedness. Thus the number
of generators is bounded by the sum of (F9.3) and $b_1(\Delta)$.
Together with (F9.2), this proves the asserted uniform bound for
every $\pi_1(\mathfrak X_N)$.

Every finite tame quotient of $\pi_1(U)$ extends to one of these
root stacks by the Abhyankar--purity argument of F8. Consequently
every finite quotient of $\pi_1^t(U;X,D)$ obeys the same bound.
The compactness argument for generating tuples now proves (F9.1).
\end{proof}

\begin{corollary}[F9b: a normal pro-$p$ kernel over a tame quotient; Proved]
Let $G$ be a quotient of $\pi_1(U)$, let $P\triangleleft G$ be
pro-$p$, and assume that the composite
\[
 \pi_1(U)\longrightarrow\Gamma:=G/P
\]
factors through $\pi_1^t(U;X,D)$. Then
\[
 d(G)\leq\max\{d_X,\lceil C_X\rceil+1\},
 \tag{F9.4}
\]
where $C_X$ is the uniform constant in F8. No solvability
hypothesis is imposed on $\Gamma$ or on its finite quotients.
\end{corollary}
\begin{proof}
Every finite simple $\mathbf F_p[G]$-module is killed by $P$,
because a normal finite $p$-group has nonzero invariants on every
nonzero representation in characteristic $p$, and the invariant
space is $G$-stable. The associated sheaf on $U$ therefore has
tame boundary monodromy. Inflation and F8 give the uniform
rank-linear bound on $H^1(G,M)$ for all these simple modules.
F9 bounds the number of generators of $\Gamma$. The relative
criterion F7a proves (F9.4).
\end{proof}

\paragraph{Remaining wild obstruction.}
The kernel of $\pi_1(U)\to\pi_1^t(U;X,D)$ is normally generated
by wild inertia groups. This does not make it a pro-$p$ group.
For example, a finite group generated by $p$-subgroups can have
minimal normal subgroups of another prime order: $S_3$ is generated
by its order-two subgroups and has normal subgroup $C_3$.
Thus F9b does not close the alternating tame/wild case or the full
local fundamental-group conjecture. No assertion of historical
priority is attached to F9 or its proof.
